\documentclass[11pt]{article}

    \usepackage[breakable]{tcolorbox}
    \usepackage{parskip} % Stop auto-indenting (to mimic markdown behaviour)
    

    % Basic figure setup, for now with no caption control since it's done
    % automatically by Pandoc (which extracts ![](path) syntax from Markdown).
    \usepackage{graphicx}
    % Keep aspect ratio if custom image width or height is specified
    \setkeys{Gin}{keepaspectratio}
    % Maintain compatibility with old templates. Remove in nbconvert 6.0
    \let\Oldincludegraphics\includegraphics
    % Ensure that by default, figures have no caption (until we provide a
    % proper Figure object with a Caption API and a way to capture that
    % in the conversion process - todo).
    \usepackage{caption}
    \DeclareCaptionFormat{nocaption}{}
    \captionsetup{format=nocaption,aboveskip=0pt,belowskip=0pt}

    \usepackage{float}
    \floatplacement{figure}{H} % forces figures to be placed at the correct location
    \usepackage{xcolor} % Allow colors to be defined
    \usepackage{enumerate} % Needed for markdown enumerations to work
    \usepackage{geometry} % Used to adjust the document margins
    \usepackage{amsmath} % Equations
    \usepackage{amssymb} % Equations
    \usepackage{textcomp} % defines textquotesingle
    % Hack from http://tex.stackexchange.com/a/47451/13684:
    \AtBeginDocument{%
        \def\PYZsq{\textquotesingle}% Upright quotes in Pygmentized code
    }
    \usepackage{upquote} % Upright quotes for verbatim code
    \usepackage{eurosym} % defines \euro

    \usepackage{iftex}
    \ifPDFTeX
        \usepackage[T1]{fontenc}
        \IfFileExists{alphabeta.sty}{
              \usepackage{alphabeta}
          }{
              \usepackage[mathletters]{ucs}
              \usepackage[utf8x]{inputenc}
          }
    \else
        \usepackage{fontspec}
        \usepackage{unicode-math}
    \fi

    \usepackage{fancyvrb} % verbatim replacement that allows latex
    \usepackage{grffile} % extends the file name processing of package graphics
                         % to support a larger range
    \makeatletter % fix for old versions of grffile with XeLaTeX
    \@ifpackagelater{grffile}{2019/11/01}
    {
      % Do nothing on new versions
    }
    {
      \def\Gread@@xetex#1{%
        \IfFileExists{"\Gin@base".bb}%
        {\Gread@eps{\Gin@base.bb}}%
        {\Gread@@xetex@aux#1}%
      }
    }
    \makeatother
    \usepackage[Export]{adjustbox} % Used to constrain images to a maximum size
    \adjustboxset{max size={0.9\linewidth}{0.9\paperheight}}

    % The hyperref package gives us a pdf with properly built
    % internal navigation ('pdf bookmarks' for the table of contents,
    % internal cross-reference links, web links for URLs, etc.)
    \usepackage{hyperref}
    % The default LaTeX title has an obnoxious amount of whitespace. By default,
    % titling removes some of it. It also provides customization options.
    \usepackage{titling}
    \usepackage{longtable} % longtable support required by pandoc >1.10
    \usepackage{booktabs}  % table support for pandoc > 1.12.2
    \usepackage{array}     % table support for pandoc >= 2.11.3
    \usepackage{calc}      % table minipage width calculation for pandoc >= 2.11.1
    \usepackage[inline]{enumitem} % IRkernel/repr support (it uses the enumerate* environment)
    \usepackage[normalem]{ulem} % ulem is needed to support strikethroughs (\sout)
                                % normalem makes italics be italics, not underlines
    \usepackage{soul}      % strikethrough (\st) support for pandoc >= 3.0.0
    \usepackage{mathrsfs}
    

    
    % Colors for the hyperref package
    \definecolor{urlcolor}{rgb}{0,.145,.698}
    \definecolor{linkcolor}{rgb}{.71,0.21,0.01}
    \definecolor{citecolor}{rgb}{.12,.54,.11}

    % ANSI colors
    \definecolor{ansi-black}{HTML}{3E424D}
    \definecolor{ansi-black-intense}{HTML}{282C36}
    \definecolor{ansi-red}{HTML}{E75C58}
    \definecolor{ansi-red-intense}{HTML}{B22B31}
    \definecolor{ansi-green}{HTML}{00A250}
    \definecolor{ansi-green-intense}{HTML}{007427}
    \definecolor{ansi-yellow}{HTML}{DDB62B}
    \definecolor{ansi-yellow-intense}{HTML}{B27D12}
    \definecolor{ansi-blue}{HTML}{208FFB}
    \definecolor{ansi-blue-intense}{HTML}{0065CA}
    \definecolor{ansi-magenta}{HTML}{D160C4}
    \definecolor{ansi-magenta-intense}{HTML}{A03196}
    \definecolor{ansi-cyan}{HTML}{60C6C8}
    \definecolor{ansi-cyan-intense}{HTML}{258F8F}
    \definecolor{ansi-white}{HTML}{C5C1B4}
    \definecolor{ansi-white-intense}{HTML}{A1A6B2}
    \definecolor{ansi-default-inverse-fg}{HTML}{FFFFFF}
    \definecolor{ansi-default-inverse-bg}{HTML}{000000}

    % common color for the border for error outputs.
    \definecolor{outerrorbackground}{HTML}{FFDFDF}

    % commands and environments needed by pandoc snippets
    % extracted from the output of `pandoc -s`
    \providecommand{\tightlist}{%
      \setlength{\itemsep}{0pt}\setlength{\parskip}{0pt}}
    \DefineVerbatimEnvironment{Highlighting}{Verbatim}{commandchars=\\\{\}}
    % Add ',fontsize=\small' for more characters per line
    \newenvironment{Shaded}{}{}
    \newcommand{\KeywordTok}[1]{\textcolor[rgb]{0.00,0.44,0.13}{\textbf{{#1}}}}
    \newcommand{\DataTypeTok}[1]{\textcolor[rgb]{0.56,0.13,0.00}{{#1}}}
    \newcommand{\DecValTok}[1]{\textcolor[rgb]{0.25,0.63,0.44}{{#1}}}
    \newcommand{\BaseNTok}[1]{\textcolor[rgb]{0.25,0.63,0.44}{{#1}}}
    \newcommand{\FloatTok}[1]{\textcolor[rgb]{0.25,0.63,0.44}{{#1}}}
    \newcommand{\CharTok}[1]{\textcolor[rgb]{0.25,0.44,0.63}{{#1}}}
    \newcommand{\StringTok}[1]{\textcolor[rgb]{0.25,0.44,0.63}{{#1}}}
    \newcommand{\CommentTok}[1]{\textcolor[rgb]{0.38,0.63,0.69}{\textit{{#1}}}}
    \newcommand{\OtherTok}[1]{\textcolor[rgb]{0.00,0.44,0.13}{{#1}}}
    \newcommand{\AlertTok}[1]{\textcolor[rgb]{1.00,0.00,0.00}{\textbf{{#1}}}}
    \newcommand{\FunctionTok}[1]{\textcolor[rgb]{0.02,0.16,0.49}{{#1}}}
    \newcommand{\RegionMarkerTok}[1]{{#1}}
    \newcommand{\ErrorTok}[1]{\textcolor[rgb]{1.00,0.00,0.00}{\textbf{{#1}}}}
    \newcommand{\NormalTok}[1]{{#1}}

    % Additional commands for more recent versions of Pandoc
    \newcommand{\ConstantTok}[1]{\textcolor[rgb]{0.53,0.00,0.00}{{#1}}}
    \newcommand{\SpecialCharTok}[1]{\textcolor[rgb]{0.25,0.44,0.63}{{#1}}}
    \newcommand{\VerbatimStringTok}[1]{\textcolor[rgb]{0.25,0.44,0.63}{{#1}}}
    \newcommand{\SpecialStringTok}[1]{\textcolor[rgb]{0.73,0.40,0.53}{{#1}}}
    \newcommand{\ImportTok}[1]{{#1}}
    \newcommand{\DocumentationTok}[1]{\textcolor[rgb]{0.73,0.13,0.13}{\textit{{#1}}}}
    \newcommand{\AnnotationTok}[1]{\textcolor[rgb]{0.38,0.63,0.69}{\textbf{\textit{{#1}}}}}
    \newcommand{\CommentVarTok}[1]{\textcolor[rgb]{0.38,0.63,0.69}{\textbf{\textit{{#1}}}}}
    \newcommand{\VariableTok}[1]{\textcolor[rgb]{0.10,0.09,0.49}{{#1}}}
    \newcommand{\ControlFlowTok}[1]{\textcolor[rgb]{0.00,0.44,0.13}{\textbf{{#1}}}}
    \newcommand{\OperatorTok}[1]{\textcolor[rgb]{0.40,0.40,0.40}{{#1}}}
    \newcommand{\BuiltInTok}[1]{{#1}}
    \newcommand{\ExtensionTok}[1]{{#1}}
    \newcommand{\PreprocessorTok}[1]{\textcolor[rgb]{0.74,0.48,0.00}{{#1}}}
    \newcommand{\AttributeTok}[1]{\textcolor[rgb]{0.49,0.56,0.16}{{#1}}}
    \newcommand{\InformationTok}[1]{\textcolor[rgb]{0.38,0.63,0.69}{\textbf{\textit{{#1}}}}}
    \newcommand{\WarningTok}[1]{\textcolor[rgb]{0.38,0.63,0.69}{\textbf{\textit{{#1}}}}}
    \makeatletter
    \newsavebox\pandoc@box
    \newcommand*\pandocbounded[1]{%
      \sbox\pandoc@box{#1}%
      % scaling factors for width and height
      \Gscale@div\@tempa\textheight{\dimexpr\ht\pandoc@box+\dp\pandoc@box\relax}%
      \Gscale@div\@tempb\linewidth{\wd\pandoc@box}%
      % select the smaller of both
      \ifdim\@tempb\p@<\@tempa\p@
        \let\@tempa\@tempb
      \fi
      % scaling accordingly (\@tempa < 1)
      \ifdim\@tempa\p@<\p@
        \scalebox{\@tempa}{\usebox\pandoc@box}%
      % scaling not needed, use as it is
      \else
        \usebox{\pandoc@box}%
      \fi
    }
    \makeatother

    % Define a nice break command that doesn't care if a line doesn't already
    % exist.
    \def\br{\hspace*{\fill} \\* }
    % Math Jax compatibility definitions
    \let\Oldtex\TeX
    \let\Oldlatex\LaTeX
    \renewcommand{\TeX}{\textrm{\Oldtex}}
    \renewcommand{\LaTeX}{\textrm{\Oldlatex}}
    % Document parameters
    % Document title

    \input{../../macros.tex}
    \title{Homework\_12}
    
    
    
    
    
    
% Pygments definitions
\makeatletter
\def\PY@reset{\let\PY@it=\relax \let\PY@bf=\relax%
    \let\PY@ul=\relax \let\PY@tc=\relax%
    \let\PY@bc=\relax \let\PY@ff=\relax}
\def\PY@tok#1{\csname PY@tok@#1\endcsname}
\def\PY@toks#1+{\ifx\relax#1\empty\else%
    \PY@tok{#1}\expandafter\PY@toks\fi}
\def\PY@do#1{\PY@bc{\PY@tc{\PY@ul{%
    \PY@it{\PY@bf{\PY@ff{#1}}}}}}}
\def\PY#1#2{\PY@reset\PY@toks#1+\relax+\PY@do{#2}}

\@namedef{PY@tok@w}{\def\PY@tc##1{\textcolor[rgb]{0.73,0.73,0.73}{##1}}}
\@namedef{PY@tok@c}{\let\PY@it=\textit\def\PY@tc##1{\textcolor[rgb]{0.24,0.48,0.48}{##1}}}
\@namedef{PY@tok@cp}{\def\PY@tc##1{\textcolor[rgb]{0.61,0.40,0.00}{##1}}}
\@namedef{PY@tok@k}{\let\PY@bf=\textbf\def\PY@tc##1{\textcolor[rgb]{0.00,0.50,0.00}{##1}}}
\@namedef{PY@tok@kp}{\def\PY@tc##1{\textcolor[rgb]{0.00,0.50,0.00}{##1}}}
\@namedef{PY@tok@kt}{\def\PY@tc##1{\textcolor[rgb]{0.69,0.00,0.25}{##1}}}
\@namedef{PY@tok@o}{\def\PY@tc##1{\textcolor[rgb]{0.40,0.40,0.40}{##1}}}
\@namedef{PY@tok@ow}{\let\PY@bf=\textbf\def\PY@tc##1{\textcolor[rgb]{0.67,0.13,1.00}{##1}}}
\@namedef{PY@tok@nb}{\def\PY@tc##1{\textcolor[rgb]{0.00,0.50,0.00}{##1}}}
\@namedef{PY@tok@nf}{\def\PY@tc##1{\textcolor[rgb]{0.00,0.00,1.00}{##1}}}
\@namedef{PY@tok@nc}{\let\PY@bf=\textbf\def\PY@tc##1{\textcolor[rgb]{0.00,0.00,1.00}{##1}}}
\@namedef{PY@tok@nn}{\let\PY@bf=\textbf\def\PY@tc##1{\textcolor[rgb]{0.00,0.00,1.00}{##1}}}
\@namedef{PY@tok@ne}{\let\PY@bf=\textbf\def\PY@tc##1{\textcolor[rgb]{0.80,0.25,0.22}{##1}}}
\@namedef{PY@tok@nv}{\def\PY@tc##1{\textcolor[rgb]{0.10,0.09,0.49}{##1}}}
\@namedef{PY@tok@no}{\def\PY@tc##1{\textcolor[rgb]{0.53,0.00,0.00}{##1}}}
\@namedef{PY@tok@nl}{\def\PY@tc##1{\textcolor[rgb]{0.46,0.46,0.00}{##1}}}
\@namedef{PY@tok@ni}{\let\PY@bf=\textbf\def\PY@tc##1{\textcolor[rgb]{0.44,0.44,0.44}{##1}}}
\@namedef{PY@tok@na}{\def\PY@tc##1{\textcolor[rgb]{0.41,0.47,0.13}{##1}}}
\@namedef{PY@tok@nt}{\let\PY@bf=\textbf\def\PY@tc##1{\textcolor[rgb]{0.00,0.50,0.00}{##1}}}
\@namedef{PY@tok@nd}{\def\PY@tc##1{\textcolor[rgb]{0.67,0.13,1.00}{##1}}}
\@namedef{PY@tok@s}{\def\PY@tc##1{\textcolor[rgb]{0.73,0.13,0.13}{##1}}}
\@namedef{PY@tok@sd}{\let\PY@it=\textit\def\PY@tc##1{\textcolor[rgb]{0.73,0.13,0.13}{##1}}}
\@namedef{PY@tok@si}{\let\PY@bf=\textbf\def\PY@tc##1{\textcolor[rgb]{0.64,0.35,0.47}{##1}}}
\@namedef{PY@tok@se}{\let\PY@bf=\textbf\def\PY@tc##1{\textcolor[rgb]{0.67,0.36,0.12}{##1}}}
\@namedef{PY@tok@sr}{\def\PY@tc##1{\textcolor[rgb]{0.64,0.35,0.47}{##1}}}
\@namedef{PY@tok@ss}{\def\PY@tc##1{\textcolor[rgb]{0.10,0.09,0.49}{##1}}}
\@namedef{PY@tok@sx}{\def\PY@tc##1{\textcolor[rgb]{0.00,0.50,0.00}{##1}}}
\@namedef{PY@tok@m}{\def\PY@tc##1{\textcolor[rgb]{0.40,0.40,0.40}{##1}}}
\@namedef{PY@tok@gh}{\let\PY@bf=\textbf\def\PY@tc##1{\textcolor[rgb]{0.00,0.00,0.50}{##1}}}
\@namedef{PY@tok@gu}{\let\PY@bf=\textbf\def\PY@tc##1{\textcolor[rgb]{0.50,0.00,0.50}{##1}}}
\@namedef{PY@tok@gd}{\def\PY@tc##1{\textcolor[rgb]{0.63,0.00,0.00}{##1}}}
\@namedef{PY@tok@gi}{\def\PY@tc##1{\textcolor[rgb]{0.00,0.52,0.00}{##1}}}
\@namedef{PY@tok@gr}{\def\PY@tc##1{\textcolor[rgb]{0.89,0.00,0.00}{##1}}}
\@namedef{PY@tok@ge}{\let\PY@it=\textit}
\@namedef{PY@tok@gs}{\let\PY@bf=\textbf}
\@namedef{PY@tok@ges}{\let\PY@bf=\textbf\let\PY@it=\textit}
\@namedef{PY@tok@gp}{\let\PY@bf=\textbf\def\PY@tc##1{\textcolor[rgb]{0.00,0.00,0.50}{##1}}}
\@namedef{PY@tok@go}{\def\PY@tc##1{\textcolor[rgb]{0.44,0.44,0.44}{##1}}}
\@namedef{PY@tok@gt}{\def\PY@tc##1{\textcolor[rgb]{0.00,0.27,0.87}{##1}}}
\@namedef{PY@tok@err}{\def\PY@bc##1{{\setlength{\fboxsep}{\string -\fboxrule}\fcolorbox[rgb]{1.00,0.00,0.00}{1,1,1}{\strut ##1}}}}
\@namedef{PY@tok@kc}{\let\PY@bf=\textbf\def\PY@tc##1{\textcolor[rgb]{0.00,0.50,0.00}{##1}}}
\@namedef{PY@tok@kd}{\let\PY@bf=\textbf\def\PY@tc##1{\textcolor[rgb]{0.00,0.50,0.00}{##1}}}
\@namedef{PY@tok@kn}{\let\PY@bf=\textbf\def\PY@tc##1{\textcolor[rgb]{0.00,0.50,0.00}{##1}}}
\@namedef{PY@tok@kr}{\let\PY@bf=\textbf\def\PY@tc##1{\textcolor[rgb]{0.00,0.50,0.00}{##1}}}
\@namedef{PY@tok@bp}{\def\PY@tc##1{\textcolor[rgb]{0.00,0.50,0.00}{##1}}}
\@namedef{PY@tok@fm}{\def\PY@tc##1{\textcolor[rgb]{0.00,0.00,1.00}{##1}}}
\@namedef{PY@tok@vc}{\def\PY@tc##1{\textcolor[rgb]{0.10,0.09,0.49}{##1}}}
\@namedef{PY@tok@vg}{\def\PY@tc##1{\textcolor[rgb]{0.10,0.09,0.49}{##1}}}
\@namedef{PY@tok@vi}{\def\PY@tc##1{\textcolor[rgb]{0.10,0.09,0.49}{##1}}}
\@namedef{PY@tok@vm}{\def\PY@tc##1{\textcolor[rgb]{0.10,0.09,0.49}{##1}}}
\@namedef{PY@tok@sa}{\def\PY@tc##1{\textcolor[rgb]{0.73,0.13,0.13}{##1}}}
\@namedef{PY@tok@sb}{\def\PY@tc##1{\textcolor[rgb]{0.73,0.13,0.13}{##1}}}
\@namedef{PY@tok@sc}{\def\PY@tc##1{\textcolor[rgb]{0.73,0.13,0.13}{##1}}}
\@namedef{PY@tok@dl}{\def\PY@tc##1{\textcolor[rgb]{0.73,0.13,0.13}{##1}}}
\@namedef{PY@tok@s2}{\def\PY@tc##1{\textcolor[rgb]{0.73,0.13,0.13}{##1}}}
\@namedef{PY@tok@sh}{\def\PY@tc##1{\textcolor[rgb]{0.73,0.13,0.13}{##1}}}
\@namedef{PY@tok@s1}{\def\PY@tc##1{\textcolor[rgb]{0.73,0.13,0.13}{##1}}}
\@namedef{PY@tok@mb}{\def\PY@tc##1{\textcolor[rgb]{0.40,0.40,0.40}{##1}}}
\@namedef{PY@tok@mf}{\def\PY@tc##1{\textcolor[rgb]{0.40,0.40,0.40}{##1}}}
\@namedef{PY@tok@mh}{\def\PY@tc##1{\textcolor[rgb]{0.40,0.40,0.40}{##1}}}
\@namedef{PY@tok@mi}{\def\PY@tc##1{\textcolor[rgb]{0.40,0.40,0.40}{##1}}}
\@namedef{PY@tok@il}{\def\PY@tc##1{\textcolor[rgb]{0.40,0.40,0.40}{##1}}}
\@namedef{PY@tok@mo}{\def\PY@tc##1{\textcolor[rgb]{0.40,0.40,0.40}{##1}}}
\@namedef{PY@tok@ch}{\let\PY@it=\textit\def\PY@tc##1{\textcolor[rgb]{0.24,0.48,0.48}{##1}}}
\@namedef{PY@tok@cm}{\let\PY@it=\textit\def\PY@tc##1{\textcolor[rgb]{0.24,0.48,0.48}{##1}}}
\@namedef{PY@tok@cpf}{\let\PY@it=\textit\def\PY@tc##1{\textcolor[rgb]{0.24,0.48,0.48}{##1}}}
\@namedef{PY@tok@c1}{\let\PY@it=\textit\def\PY@tc##1{\textcolor[rgb]{0.24,0.48,0.48}{##1}}}
\@namedef{PY@tok@cs}{\let\PY@it=\textit\def\PY@tc##1{\textcolor[rgb]{0.24,0.48,0.48}{##1}}}

\def\PYZbs{\char`\\}
\def\PYZus{\char`\_}
\def\PYZob{\char`\{}
\def\PYZcb{\char`\}}
\def\PYZca{\char`\^}
\def\PYZam{\char`\&}
\def\PYZlt{\char`\<}
\def\PYZgt{\char`\>}
\def\PYZsh{\char`\#}
\def\PYZpc{\char`\%}
\def\PYZdl{\char`\$}
\def\PYZhy{\char`\-}
\def\PYZsq{\char`\'}
\def\PYZdq{\char`\"}
\def\PYZti{\char`\~}
% for compatibility with earlier versions
\def\PYZat{@}
\def\PYZlb{[}
\def\PYZrb{]}
\makeatother


    % For linebreaks inside Verbatim environment from package fancyvrb.
    \makeatletter
        \newbox\Wrappedcontinuationbox
        \newbox\Wrappedvisiblespacebox
        \newcommand*\Wrappedvisiblespace {\textcolor{red}{\textvisiblespace}}
        \newcommand*\Wrappedcontinuationsymbol {\textcolor{red}{\llap{\tiny$\m@th\hookrightarrow$}}}
        \newcommand*\Wrappedcontinuationindent {3ex }
        \newcommand*\Wrappedafterbreak {\kern\Wrappedcontinuationindent\copy\Wrappedcontinuationbox}
        % Take advantage of the already applied Pygments mark-up to insert
        % potential linebreaks for TeX processing.
        %        {, <, #, %, $, ' and ": go to next line.
        %        _, }, ^, &, >, - and ~: stay at end of broken line.
        % Use of \textquotesingle for straight quote.
        \newcommand*\Wrappedbreaksatspecials {%
            \def\PYGZus{\discretionary{\char`\_}{\Wrappedafterbreak}{\char`\_}}%
            \def\PYGZob{\discretionary{}{\Wrappedafterbreak\char`\{}{\char`\{}}%
            \def\PYGZcb{\discretionary{\char`\}}{\Wrappedafterbreak}{\char`\}}}%
            \def\PYGZca{\discretionary{\char`\^}{\Wrappedafterbreak}{\char`\^}}%
            \def\PYGZam{\discretionary{\char`\&}{\Wrappedafterbreak}{\char`\&}}%
            \def\PYGZlt{\discretionary{}{\Wrappedafterbreak\char`\<}{\char`\<}}%
            \def\PYGZgt{\discretionary{\char`\>}{\Wrappedafterbreak}{\char`\>}}%
            \def\PYGZsh{\discretionary{}{\Wrappedafterbreak\char`\#}{\char`\#}}%
            \def\PYGZpc{\discretionary{}{\Wrappedafterbreak\char`\%}{\char`\%}}%
            \def\PYGZdl{\discretionary{}{\Wrappedafterbreak\char`\$}{\char`\$}}%
            \def\PYGZhy{\discretionary{\char`\-}{\Wrappedafterbreak}{\char`\-}}%
            \def\PYGZsq{\discretionary{}{\Wrappedafterbreak\textquotesingle}{\textquotesingle}}%
            \def\PYGZdq{\discretionary{}{\Wrappedafterbreak\char`\"}{\char`\"}}%
            \def\PYGZti{\discretionary{\char`\~}{\Wrappedafterbreak}{\char`\~}}%
        }
        % Some characters . , ; ? ! / are not pygmentized.
        % This macro makes them "active" and they will insert potential linebreaks
        \newcommand*\Wrappedbreaksatpunct {%
            \lccode`\~`\.\lowercase{\def~}{\discretionary{\hbox{\char`\.}}{\Wrappedafterbreak}{\hbox{\char`\.}}}%
            \lccode`\~`\,\lowercase{\def~}{\discretionary{\hbox{\char`\,}}{\Wrappedafterbreak}{\hbox{\char`\,}}}%
            \lccode`\~`\;\lowercase{\def~}{\discretionary{\hbox{\char`\;}}{\Wrappedafterbreak}{\hbox{\char`\;}}}%
            \lccode`\~`\:\lowercase{\def~}{\discretionary{\hbox{\char`\:}}{\Wrappedafterbreak}{\hbox{\char`\:}}}%
            \lccode`\~`\?\lowercase{\def~}{\discretionary{\hbox{\char`\?}}{\Wrappedafterbreak}{\hbox{\char`\?}}}%
            \lccode`\~`\!\lowercase{\def~}{\discretionary{\hbox{\char`\!}}{\Wrappedafterbreak}{\hbox{\char`\!}}}%
            \lccode`\~`\/\lowercase{\def~}{\discretionary{\hbox{\char`\/}}{\Wrappedafterbreak}{\hbox{\char`\/}}}%
            \catcode`\.\active
            \catcode`\,\active
            \catcode`\;\active
            \catcode`\:\active
            \catcode`\?\active
            \catcode`\!\active
            \catcode`\/\active
            \lccode`\~`\~
        }
    \makeatother

    \let\OriginalVerbatim=\Verbatim
    \makeatletter
    \renewcommand{\Verbatim}[1][1]{%
        %\parskip\z@skip
        \sbox\Wrappedcontinuationbox {\Wrappedcontinuationsymbol}%
        \sbox\Wrappedvisiblespacebox {\FV@SetupFont\Wrappedvisiblespace}%
        \def\FancyVerbFormatLine ##1{\hsize\linewidth
            \vtop{\raggedright\hyphenpenalty\z@\exhyphenpenalty\z@
                \doublehyphendemerits\z@\finalhyphendemerits\z@
                \strut ##1\strut}%
        }%
        % If the linebreak is at a space, the latter will be displayed as visible
        % space at end of first line, and a continuation symbol starts next line.
        % Stretch/shrink are however usually zero for typewriter font.
        \def\FV@Space {%
            \nobreak\hskip\z@ plus\fontdimen3\font minus\fontdimen4\font
            \discretionary{\copy\Wrappedvisiblespacebox}{\Wrappedafterbreak}
            {\kern\fontdimen2\font}%
        }%

        % Allow breaks at special characters using \PYG... macros.
        \Wrappedbreaksatspecials
        % Breaks at punctuation characters . , ; ? ! and / need catcode=\active
        \OriginalVerbatim[#1,codes*=\Wrappedbreaksatpunct]%
    }
    \makeatother

    % Exact colors from NB
    \definecolor{incolor}{HTML}{303F9F}
    \definecolor{outcolor}{HTML}{D84315}
    \definecolor{cellborder}{HTML}{CFCFCF}
    \definecolor{cellbackground}{HTML}{F7F7F7}

    % prompt
    \makeatletter
    \newcommand{\boxspacing}{\kern\kvtcb@left@rule\kern\kvtcb@boxsep}
    \makeatother
    \newcommand{\prompt}[4]{
        {\ttfamily\llap{{\color{#2}[#3]:\hspace{3pt}#4}}\vspace{-\baselineskip}}
    }
    

    
    % Prevent overflowing lines due to hard-to-break entities
    \sloppy
    % Setup hyperref package
    \hypersetup{
      breaklinks=true,  % so long urls are correctly broken across lines
      colorlinks=true,
      urlcolor=urlcolor,
      linkcolor=linkcolor,
      citecolor=citecolor,
      }
    % Slightly bigger margins than the latex defaults
    
    \geometry{verbose,tmargin=1in,bmargin=1in,lmargin=1in,rmargin=1in}
    
    

\begin{document}
    
    \maketitle
    
    

    
    Probability for Data Science

UC Berkeley, Spring 2026

Ani Adhikari

CC BY-NC-SA 4.0

This content is protected and may not be shared, uploaded, or
distributed.

    \hypertarget{homework-12-due-monday-april-27th-at-5-pm}{%
\section{Homework 12 (Due Monday, April 27th at 5
PM)}\label{homework-12-due-monday-april-27th-at-5-pm}}

    \begin{tcolorbox}[breakable, size=fbox, boxrule=1pt, pad at break*=1mm,colback=cellbackground, colframe=cellborder]
\prompt{In}{incolor}{ }{\boxspacing}
\begin{Verbatim}[commandchars=\\\{\}]
\PY{k+kn}{import}\PY{+w}{ }\PY{n+nn}{warnings}
\PY{n}{warnings}\PY{o}{.}\PY{n}{filterwarnings}\PY{p}{(}\PY{l+s+s1}{\PYZsq{}}\PY{l+s+s1}{ignore}\PY{l+s+s1}{\PYZsq{}}\PY{p}{)}

\PY{k+kn}{from}\PY{+w}{ }\PY{n+nn}{prob140}\PY{+w}{ }\PY{k+kn}{import} \PY{o}{*}
\PY{k+kn}{from}\PY{+w}{ }\PY{n+nn}{datascience}\PY{+w}{ }\PY{k+kn}{import} \PY{o}{*}
\PY{k+kn}{import}\PY{+w}{ }\PY{n+nn}{numpy}\PY{+w}{ }\PY{k}{as}\PY{+w}{ }\PY{n+nn}{np}
\PY{k+kn}{from}\PY{+w}{ }\PY{n+nn}{scipy}\PY{+w}{ }\PY{k+kn}{import} \PY{n}{stats}

\PY{k+kn}{import}\PY{+w}{ }\PY{n+nn}{matplotlib}\PY{n+nn}{.}\PY{n+nn}{pyplot}\PY{+w}{ }\PY{k}{as}\PY{+w}{ }\PY{n+nn}{plt}
\PY{o}{\PYZpc{}}\PY{k}{matplotlib} inline
\PY{k+kn}{import}\PY{+w}{ }\PY{n+nn}{matplotlib}
\PY{n}{matplotlib}\PY{o}{.}\PY{n}{style}\PY{o}{.}\PY{n}{use}\PY{p}{(}\PY{l+s+s1}{\PYZsq{}}\PY{l+s+s1}{fivethirtyeight}\PY{l+s+s1}{\PYZsq{}}\PY{p}{)}
\end{Verbatim}
\end{tcolorbox}

    \begin{tcolorbox}[breakable, size=fbox, boxrule=1pt, pad at break*=1mm,colback=cellbackground, colframe=cellborder]
\prompt{In}{incolor}{ }{\boxspacing}
\begin{Verbatim}[commandchars=\\\{\}]
\PY{k+kn}{from}\PY{+w}{ }\PY{n+nn}{hw12\PYZus{}grader}\PY{+w}{ }\PY{k+kn}{import} \PY{n}{Autograder}
\PY{n}{grader} \PY{o}{=} \PY{n}{Autograder}\PY{p}{(}\PY{p}{)}
\end{Verbatim}
\end{tcolorbox}

    \hypertarget{instructions}{%
\subsubsection{Instructions}\label{instructions}}

\hypertarget{questions-1-2-and-3-can-be-answered-based-on-tuesdays-lecture-and-past-content-alone.}{%
\paragraph{\texorpdfstring{\textbf{Questions 1, 2, and 3 can be answered
based on Tuesday's lecture (and past content)
alone.}}{Questions 1, 2, and 3 can be answered based on Tuesday's lecture (and past content) alone.}}\label{questions-1-2-and-3-can-be-answered-based-on-tuesdays-lecture-and-past-content-alone.}}

\hypertarget{we-strongly-encourage-you-to-attempt-questions-1-and-2-independently-relying-on-the-provided-textbook-sections-your-own-reasoning-and-our-ai-tools-see-below.-if-you-find-yourself-stuck-we-recommend-attending-office-hours-or-homework-party-or-posting-your-questions-on-ed.-if-youre-having-trouble-getting-started-on-problems-without-the-use-of-ai-we-encourage-you-to-follow-our-recommendation.}{%
\paragraph{\texorpdfstring{\textbf{We strongly encourage you to attempt
Questions 1 and 2 independently, relying on the provided textbook
sections, your own reasoning, and our AI tools (see below). If you find
yourself stuck, we recommend attending office hours or homework party or
posting your questions on Ed. If you're having trouble getting started
on problems without the use of AI, we encourage you to follow our
recommendation.}}{We strongly encourage you to attempt Questions 1 and 2 independently, relying on the provided textbook sections, your own reasoning, and our AI tools (see below). If you find yourself stuck, we recommend attending office hours or homework party or posting your questions on Ed. If you're having trouble getting started on problems without the use of AI, we encourage you to follow our recommendation.}}\label{we-strongly-encourage-you-to-attempt-questions-1-and-2-independently-relying-on-the-provided-textbook-sections-your-own-reasoning-and-our-ai-tools-see-below.-if-you-find-yourself-stuck-we-recommend-attending-office-hours-or-homework-party-or-posting-your-questions-on-ed.-if-youre-having-trouble-getting-started-on-problems-without-the-use-of-ai-we-encourage-you-to-follow-our-recommendation.}}

Your homeworks will generally have two components: a written portion and
a portion that also involves code. Written work should be completed on
paper, and coding questions should be done in the notebook. Start the
work for the written portions of each section on a new page. You are
welcome to \(\LaTeX\) your answers to the written portions, but staff
will not be able to assist you with \(\LaTeX\) related issues.

It is your responsibility to ensure that both components of the homework
are submitted completely and properly to Gradescope. \textbf{Make sure
to assign each page of your pdf to the correct question. Refer to the
bottom of the notebook for submission instructions.}

    \hypertarget{how-to-do-your-homework}{%
\subsubsection{How to Do Your Homework}\label{how-to-do-your-homework}}

The point of homework is for you to try your hand at using what you've
learned in class. The steps to follow:

\begin{itemize}
\tightlist
\item
  Go to lecture and sections, and also go over the relevant text
  sections before starting on the homework. This will remind you what
  was covered in class, and the text will typically contain examples not
  covered in lecture. The weekly Study Guide will list what you should
  read.
\item
  Work on some of the practice problems before starting on the homework.
\item
  Attempt the homework problems by yourself with the text, section work,
  and practice materials all at hand. Sometimes the week's lab will help
  as well. The two steps above will help this step go faster and be more
  fruitful.
\item
  At this point, seek help if you need it. Don't ask how to do the
  problem --- ask how to get started, or for a nudge to get you past
  where you are stuck. Always say what you have already tried. That
  helps us help you more effectively.
\item
  For a good measure of your understanding, keep track of the fraction
  of the homework you can do by yourself or with minimal help. It's a
  better measure than your homework score, and only you can measure it.
\end{itemize}

    \hypertarget{rules-for-homework}{%
\subsubsection{Rules for Homework}\label{rules-for-homework}}

\begin{itemize}
\tightlist
\item
  Every answer should contain a calculation or reasoning. For example, a
  calculation such as \((1/3)(0.8) + (2/3)(0.7)\) or
  \texttt{sum({[}(1/3)*0.8,\ (2/3)*0.7{]})}is fine without further
  explanation or simplification. If we want you to simplify, we'll ask
  you to. But just \({5 \choose 2}\) by itself is not fine; write ``we
  want any 2 out of the 5 frogs and they can appear in any order'' or
  whatever reasoning you used. Reasoning can be brief and abbreviated,
  e.g.~``product rule'' or ``not mutually exclusive.''
\item
  You may consult others (see ``How to Do Your Homework'' above) but you
  must write up your own answers using your own words, notation, and
  sequence of steps.
\item
  We'll be using Gradescope. You must submit the homework according to
  the instructions at the end of homework set.
\end{itemize}

\hypertarget{must-read-special-rules-on-this-hw-for-jupytutor}{%
\subsubsection{\texorpdfstring{\textbf{{[}MUST READ{]} Special Rules on
this HW for
Jupytutor}}{{[}MUST READ{]} Special Rules on this HW for Jupytutor}}\label{must-read-special-rules-on-this-hw-for-jupytutor}}

\begin{itemize}
\tightlist
\item
  In this homework, we are demoing a new AI Tutor for Data 140 called
  Jupytutor.
\item
  The tutor will be available to help you solve Question 2 only. It will
  appear the below autograder cell if you fail the test case.
\item
  If you do not want to use the tutor, click the 3 dots in the bottom
  left corner, and click ``Turn off Jupytutor for this notebook''.
\item
  It should not appear below any other cells, but if it does, you should
  ignore the outputs of those windows as the AI is missing context for
  the other questions.
\item
  \textbf{Grading for Question 2:} Although we provide an autograder for
  Q2 to check your work, you must still show your work and answer in the
  written part of the homework. \textbf{Passing the autograder alone
  does not guarantee full credit for Q2}.
\item
  Recommended Usage: ask ``how do I start/continue this question?'' if
  you are stuck. The AI is familiar with parts of our course textbook
  and notation.
\item
  If any technical issues with Jupytutor arise, please post on Ed.
\end{itemize}

\hypertarget{we-will-not-grade-assignments-which-do-not-have-pages-correctly-selected-for-each-question.}{%
\subsection{\texorpdfstring{\textbf{We will not grade assignments which
do not have pages correctly selected for each
question.}}{We will not grade assignments which do not have pages correctly selected for each question.}}\label{we-will-not-grade-assignments-which-do-not-have-pages-correctly-selected-for-each-question.}}

    \newpage

    \hypertarget{overlapping-counts}{%
\subsection{1. Overlapping Counts}\label{overlapping-counts}}

Consider a sequence of i.i.d. Bernoulli \((p)\) trials. Consider the
three variables \(X\), \(Y\), and \(V\) defined by:

\begin{itemize}
\tightlist
\item
  \(X\) is the number of successes in trials 1 through 100
\item
  \(Y\) is the number of successes in trials 51 through 100
\item
  \(V\) is the number of successes in trials 51 through 150
\end{itemize}

\textbf{a)} For each of \(X\), \(Y\), and \(V\), say what the
distribution is and provide the parameters.


$X \sim \Binom(100, p), Y \sim \Binom(50, p), V \sim \Binom(100, p) $ 



\textbf{b)} Fix \(k\) in the range \(0, 1, \ldots, 100\) and find the
conditional distribution of \(Y\) given \(X = k\). Recognize this as a
famous one and provide the parameters. Start with \(P(Y = y|X=x)\) and
use the
\href{https://data140.org/textbook/content/chapter-02/multiplication/\#division-rule}{Division
Rule}.


Let $Z$  be the number of successes in trials 1 through 50, $\sim \Binom(50, p)$ 

\begin{align*}
	P(Y=y \mid X = x) &= \frac{P(Y=y, X=x)}{P(X=x)} \\
&= \frac{P(Z = x-y)P(Y=y)}{P(X=x)} \\
\end{align*}

\begin{align*}
	P(y=y \mid X=k) &= \frac{P(Z = k-y) P(Y=y)}{P(X=k)} \\
	&= \binom{50}{k-y} p^{k-y}(1-p)^{50-k+y} \binom{50}{y}p^{y} (1-p)^{50-y} / \binom{100}{k}p^{k} (1-p)^{100-k}\\
	&= \binom{50}{k-y} \binom{50}{y} p^{k} (1-p)^{100-k} / \binom{100}{k}p^{k} (1-p)^{100-k}\\
	&= \binom{50}{k-y} \binom{50}{y} / \binom{100}{k} \\
\end{align*}

This is the hypergeometric(100, 50, k)  distribution.


\textbf{c)} Find the least squares predictor of \(Y\) based on \(X\) and
say whether it is a linear function of \(X\). (If it is, then the best
linear predictor is in fact the best among all predictors.) Find
\(Var(Y \mid X)\). For a reminder of the least squares predictor, read
through
\href{https://data140.org/textbook/content/chapter-22/least-squares-predictor/\#best-predictor}{Ch.
22.2.2} (Least Squares Predictor).


\begin{align*}
	E[Y \mid X] &= X \frac{50}{100} \\
	&= \frac{1}{2}X \\
\end{align*}

This is a linear function of $X$. 


\begin{align*}
	\Var(Y \mid X) &= X \frac{50}{100} \cdot \frac{50}{100} \cdot \frac{100-X}{100-1} \\
	&= \frac{1}{4} X \cdot \frac{100-X}{99} \\
\end{align*}



\textbf{d)} Find \(E(V \mid X)\), \(Var(V \mid X)\), and the correlation
\(r(X, V)\). The equation for correlation can be found
\href{https://data140.org/textbook/content/chapter-13/covariance/\#correlation}{here}.

Let $Z$  be the number of successes in trials between 101 and 150.


\begin{align*}
	E[V \mid X] &= E[Y + Z \mid X] \\
	&= E[Y \mid X] + E[Z \mid X]	\\
	&= E[Y \mid X]  + E[Z] \\
	&= \frac{1}{2}X + 50p \\
\end{align*}


\begin{align*}
	\Var(V \mid X) &= \Var(Y + Z \mid X) \\
	&= \Var(Y \mid X) + \Var(Z \mid X) \\
	&= \frac{1}{4} X \cdot \frac{100-X}{99} + 50p(1-p) \\
\end{align*}


$$
	r(X,V) = \frac{\Cov(X,V)}{SD(X)SD(V)} 
$$


\begin{align*}
	\Cov(X, V) &= \Cov(X, Y + Z) \\
	&= \Cov(X, Y) + \Cov(X, Z) \\
	&= \Cov(X, Y) \\
\end{align*}

Let $I_{i}$ be an indicator variable indicating if trial $i$ is a success.

\begin{align*}
	\Cov(X,Y) &= \Cov(\sum_{i=1}^{100} I_{i}, \sum_{i=51}^{100} I_{i})\\
&= \Cov(\sum_{i=1}^{50} I_{i}, 0) + \Cov(\sum_{i=51}^{100} I_{i}, \sum_{i=51}^{100} I_{i}) \\
&= \Var(\sum_{i=51}^{100} I_{i}) \\
&= 50 \Var(I_{1}) \\
&= 50 p(1-p) \\
\end{align*}


\begin{align*}
	\Cov(X, Y) &= 50p(1-p) \\
	r(X,V) &= \frac{50p(1-p)}{\sqrt{100 p(1-p)}{\sqrt{100 p(1-p)}}} \\
	&= \frac{50p(1-p)}{100p(1-p)} \\
	&= \frac{1}{2} \\
\end{align*}


\textbf{e)} Simulate 20,000 \((X, V)\) pairs for \(p = 0.5\) and draw
the scatter plot of the observed points. Plot \(E(V \mid X)\) as a
function of \(X\) on the same plot. Use the cell below. The arrays
\texttt{x} and \texttt{v} should contain the observed values of \(X\)
and \(V\). The array \texttt{exp\_V\_given\_x} should contain
\(E(V \mid X = x)\) for each \(x\) in \texttt{x}, using the formula you
derived in \textbf{d}.
 
    \begin{tcolorbox}[breakable, size=fbox, boxrule=1pt, pad at break*=1mm,colback=cellbackground, colframe=cellborder]
\prompt{In}{incolor}{25}{\boxspacing}
\begin{Verbatim}[commandchars=\\\{\}]
\PY{c+c1}{\PYZsh{} Simulation for e}

\PY{n}{p} \PY{o}{=} \PY{l+m+mf}{0.5}

\PY{n}{y} \PY{o}{=} \PY{n}{stats}\PY{o}{.}\PY{n}{binom}\PY{o}{.}\PY{n}{rvs}\PY{p}{(}\PY{l+m+mi}{50}\PY{p}{,} \PY{n}{p}\PY{p}{,} \PY{n}{size} \PY{o}{=} \PY{l+m+mi}{20\PYZus{}000}\PY{p}{)}
\PY{n}{x} \PY{o}{=} \PY{n}{stats}\PY{o}{.}\PY{n}{binom}\PY{o}{.}\PY{n}{rvs}\PY{p}{(}\PY{l+m+mi}{50}\PY{p}{,} \PY{n}{p}\PY{p}{,} \PY{n}{size} \PY{o}{=} \PY{l+m+mi}{20\PYZus{}000}\PY{p}{)} \PY{o}{+} \PY{n}{y}
\PY{n}{v} \PY{o}{=} \PY{n}{y} \PY{o}{+} \PY{n}{stats}\PY{o}{.}\PY{n}{binom}\PY{o}{.}\PY{n}{rvs}\PY{p}{(}\PY{l+m+mi}{50}\PY{p}{,} \PY{n}{p}\PY{p}{,} \PY{n}{size}\PY{o}{=} \PY{l+m+mi}{20\PYZus{}000}\PY{p}{)}
\PY{n}{exp\PYZus{}V\PYZus{}given\PYZus{}x} \PY{o}{=} \PY{p}{(}\PY{l+m+mf}{0.5} \PY{o}{*} \PY{n}{x}\PY{p}{)} \PY{o}{+} \PY{l+m+mi}{50} \PY{o}{*} \PY{n}{p}

\PY{c+c1}{\PYZsh{} Don\PYZsq{}t change the lines below}
\PY{n}{plt}\PY{o}{.}\PY{n}{figure}\PY{p}{(}\PY{n}{figsize}\PY{o}{=}\PY{p}{(}\PY{l+m+mi}{6}\PY{p}{,} \PY{l+m+mi}{6}\PY{p}{)}\PY{p}{)}
\PY{n}{plt}\PY{o}{.}\PY{n}{axes}\PY{p}{(}\PY{p}{)}\PY{o}{.}\PY{n}{set\PYZus{}aspect}\PY{p}{(}\PY{l+s+s1}{\PYZsq{}}\PY{l+s+s1}{equal}\PY{l+s+s1}{\PYZsq{}}\PY{p}{)}
\PY{n}{plt}\PY{o}{.}\PY{n}{xticks}\PY{p}{(}\PY{n}{np}\PY{o}{.}\PY{n}{arange}\PY{p}{(}\PY{l+m+mi}{30}\PY{p}{,} \PY{l+m+mi}{71}\PY{p}{,} \PY{l+m+mi}{10}\PY{p}{)}\PY{p}{)}
\PY{n}{plt}\PY{o}{.}\PY{n}{yticks}\PY{p}{(}\PY{n}{np}\PY{o}{.}\PY{n}{arange}\PY{p}{(}\PY{l+m+mi}{30}\PY{p}{,} \PY{l+m+mi}{71}\PY{p}{,} \PY{l+m+mi}{10}\PY{p}{)}\PY{p}{)}
\PY{n}{plt}\PY{o}{.}\PY{n}{scatter}\PY{p}{(}\PY{n}{x}\PY{p}{,} \PY{n}{v}\PY{p}{,} \PY{n}{color}\PY{o}{=}\PY{l+s+s1}{\PYZsq{}}\PY{l+s+s1}{darkblue}\PY{l+s+s1}{\PYZsq{}}\PY{p}{,} \PY{n}{s}\PY{o}{=}\PY{l+m+mi}{10}\PY{p}{)}
\PY{n}{plt}\PY{o}{.}\PY{n}{scatter}\PY{p}{(}\PY{n}{x}\PY{p}{,} \PY{n}{exp\PYZus{}V\PYZus{}given\PYZus{}x}\PY{p}{,} \PY{n}{color}\PY{o}{=}\PY{l+s+s1}{\PYZsq{}}\PY{l+s+s1}{gold}\PY{l+s+s1}{\PYZsq{}}\PY{p}{,} \PY{n}{s}\PY{o}{=}\PY{l+m+mi}{10}\PY{p}{)}\PY{p}{;}
\end{Verbatim}
\end{tcolorbox}

    \begin{center}
    \adjustimage{max size={0.9\linewidth}{0.9\paperheight}}{output_9_0.png}
    \end{center}
    { \hspace*{\fill} \\}
    
    Complete the cell below so that the \texttt{x\_su} is \texttt{x} in
standard units, \texttt{v\_su} is \texttt{v} in standard units, and the
last line evaluates to the observed correlation between \texttt{x} and
\texttt{v}. Check that the output is consistent with your calculation in
\textbf{d}.

    \begin{tcolorbox}[breakable, size=fbox, boxrule=1pt, pad at break*=1mm,colback=cellbackground, colframe=cellborder]
\prompt{In}{incolor}{26}{\boxspacing}
\begin{Verbatim}[commandchars=\\\{\}]
\PY{c+c1}{\PYZsh{} Part e continued}

\PY{n}{x\PYZus{}su} \PY{o}{=} \PY{p}{(}\PY{n}{x} \PY{o}{\PYZhy{}} \PY{n}{np}\PY{o}{.}\PY{n}{mean}\PY{p}{(}\PY{n}{x}\PY{p}{)}\PY{p}{)} \PY{o}{/} \PY{n}{np}\PY{o}{.}\PY{n}{std}\PY{p}{(}\PY{n}{x}\PY{p}{)}
\PY{n}{v\PYZus{}su} \PY{o}{=} \PY{p}{(}\PY{n}{v} \PY{o}{\PYZhy{}} \PY{n}{np}\PY{o}{.}\PY{n}{mean}\PY{p}{(}\PY{n}{v}\PY{p}{)}\PY{p}{)} \PY{o}{/} \PY{n}{np}\PY{o}{.}\PY{n}{std}\PY{p}{(}\PY{n}{v}\PY{p}{)}

\PY{c+c1}{\PYZsh{} expectation of x * v }

\PY{l+m+mi}{1}\PY{o}{/}\PY{l+m+mi}{20\PYZus{}000} \PY{o}{*} \PY{n}{np}\PY{o}{.}\PY{n}{sum}\PY{p}{(}\PY{n}{x\PYZus{}su} \PY{o}{*} \PY{n}{v\PYZus{}su}\PY{p}{)}
\end{Verbatim}
\end{tcolorbox}

            \begin{tcolorbox}[breakable, size=fbox, boxrule=.5pt, pad at break*=1mm, opacityfill=0]
\prompt{Out}{outcolor}{26}{\boxspacing}
\begin{Verbatim}[commandchars=\\\{\}]
0.50436291167647651
\end{Verbatim}
\end{tcolorbox}
          \newpage

    \hypertarget{midterms-1-and-2}{%
\subsection{2. Midterms 1 and 2}\label{midterms-1-and-2}}

Before attempting this problem, read through
\href{https://data140.org/textbook/content/chapter-23/multivariate-normal-vectors/\#sum-and-difference}{Ch.
23.2.2} (Sum and Difference of a Bivariate Normal). For a problem
related to this one, look at Ch. 23.5 Ex.3, which will be done in
Wednesday's section.

In a large class, the distribution of scores in Midterms 1 and 2 is
approximately bivariate normal. Define the following notation.

\begin{itemize}
\tightlist
\item
  For \(i = 1, 2\), \(\mu_i\) is the class average on Midterm \(i\), and
  \(\sigma_i^2\) is the class variance on Midterm \(i\)
\item
  The correlation between Midterm 1 and 2 scores in the class is
  \(\rho\)
\end{itemize}

For both questions below, \textbf{find the answer in terms of rho or
constants}. Then, set the variables in the cells below to the exact
answer in the case that \(\rho\) = 0.6, and double check your solution
with our autograder.

    \textbf{a)} For approximately what proportion of the students was the
Midterm 2 score (in standard units) higher than the Midterm 1 score (in
standard units)? First \textbf{find the answer in terms of rho or
constants}, then find the answer numerically in the cell below.
\begin{align*}
P(Z_2 \ge Z_1) &= P(Z_2 - Z_1 \ge 0) \\
\end{align*}

$D = Z_2 - Z_1 \sim$ Normal($0, 2-2\rho = 2(1-\rho)$)
\begin{align*}
	P(Z_2 - Z_1 \ge 0) &= P(\frac{(D) - (0)}{2(1-\rho)} \ge 0)\\
			   &= \Phi(0)
\end{align*}

    \begin{tcolorbox}[breakable, size=fbox, boxrule=1pt, pad at break*=1mm,colback=cellbackground, colframe=cellborder]
\prompt{In}{incolor}{8}{\boxspacing}
\begin{Verbatim}[commandchars=\\\{\}]
\PY{c+c1}{\PYZsh{} answer to a}
\PY{n}{rho} \PY{o}{=} \PY{l+m+mf}{0.6} \PY{c+c1}{\PYZsh{} don\PYZsq{}t change this line}
\PY{n}{q2a\PYZus{}answer} \PY{o}{=} \PY{l+m+mf}{.5}
\PY{n}{q2a\PYZus{}answer}
\end{Verbatim}
\end{tcolorbox}

            \begin{tcolorbox}[breakable, size=fbox, boxrule=.5pt, pad at break*=1mm, opacityfill=0]
\prompt{Out}{outcolor}{8}{\boxspacing}
\begin{Verbatim}[commandchars=\\\{\}]
0.5
\end{Verbatim}
\end{tcolorbox}
        
    \begin{tcolorbox}[breakable, size=fbox, boxrule=1pt, pad at break*=1mm,colback=cellbackground, colframe=cellborder]
\prompt{In}{incolor}{9}{\boxspacing}
\begin{Verbatim}[commandchars=\\\{\}]
\PY{n}{grader}\PY{o}{.}\PY{n}{grade\PYZus{}q2a}\PY{p}{(}\PY{n}{q2a\PYZus{}answer}\PY{p}{)} \PY{c+c1}{\PYZsh{} don\PYZsq{}t change this line}
\end{Verbatim}
\end{tcolorbox}

    \begin{Verbatim}[commandchars=\\\{\}]
CORRECT
    \end{Verbatim}

    \textbf{b)} For approximately what proportion of the students did the
two midterm scores (each in standard units) differ by more than 1? Round
your answer to the nearest 5 decimal places. First \textbf{find the
answer in terms of rho or constants}, then find the answer numerically
in the cell below.

\begin{align*}
	P(\vert D \vert  \ge 1) &= 1 - P(-1 \le D \le 1)  \\
	&= 1 - (P(D \le 1) - P(D \le -1)) \\
	&= 1 - (\Phi(\frac{1}{\sqrt{2(1-\rho)}}) - \Phi(-\frac{1}{\sqrt{2(1-\rho)}})) \\
\end{align*}


    \begin{tcolorbox}[breakable, size=fbox, boxrule=1pt, pad at break*=1mm,colback=cellbackground, colframe=cellborder]
\prompt{In}{incolor}{21}{\boxspacing}
\begin{Verbatim}[commandchars=\\\{\}]
\PY{c+c1}{\PYZsh{} answer to b}
\PY{n}{rho} \PY{o}{=} \PY{l+m+mf}{0.6} \PY{c+c1}{\PYZsh{} don\PYZsq{}t change this line}
\PY{n}{q2b\PYZus{}answer} \PY{o}{=} \PY{l+m+mi}{1} \PY{o}{\PYZhy{}} \PY{p}{(}\PY{n}{stats}\PY{o}{.}\PY{n}{norm}\PY{o}{.}\PY{n}{cdf}\PY{p}{(}\PY{l+m+mi}{1}\PY{o}{/}\PY{n}{np}\PY{o}{.}\PY{n}{sqrt}\PY{p}{(}\PY{l+m+mi}{2}\PY{o}{*}\PY{p}{(}\PY{l+m+mi}{1}\PY{o}{\PYZhy{}}\PY{l+m+mf}{0.6}\PY{p}{)}\PY{p}{)}\PY{p}{)} \PY{o}{\PYZhy{}} \PY{n}{stats}\PY{o}{.}\PY{n}{norm}\PY{o}{.}\PY{n}{cdf}\PY{p}{(}\PY{o}{\PYZhy{}}\PY{l+m+mi}{1}\PY{o}{/}\PY{n}{np}\PY{o}{.}\PY{n}{sqrt}\PY{p}{(}\PY{l+m+mi}{2}\PY{o}{*}\PY{p}{(}\PY{l+m+mi}{1}\PY{o}{\PYZhy{}}\PY{l+m+mf}{0.6}\PY{p}{)}\PY{p}{)}\PY{p}{)}\PY{p}{)}

\PY{n}{q2b\PYZus{}answer} \PY{o}{=} \PY{n}{np}\PY{o}{.}\PY{n}{round}\PY{p}{(}\PY{n}{q2b\PYZus{}answer}\PY{p}{,} \PY{l+m+mi}{5}\PY{p}{)} \PY{c+c1}{\PYZsh{} don\PYZsq{}t change this line, this rounds your answer}
\PY{n+nb}{print}\PY{p}{(}\PY{l+s+sa}{f}\PY{l+s+s2}{\PYZdq{}}\PY{l+s+si}{\PYZob{}}\PY{n}{q2b\PYZus{}answer}\PY{l+s+si}{:}\PY{l+s+s2}{.5f}\PY{l+s+si}{\PYZcb{}}\PY{l+s+s2}{\PYZdq{}}\PY{p}{)}
\end{Verbatim}
\end{tcolorbox}

    \begin{Verbatim}[commandchars=\\\{\}]
0.26355
    \end{Verbatim}

    \begin{tcolorbox}[breakable, size=fbox, boxrule=1pt, pad at break*=1mm,colback=cellbackground, colframe=cellborder]
\prompt{In}{incolor}{22}{\boxspacing}
\begin{Verbatim}[commandchars=\\\{\}]
\PY{n}{grader}\PY{o}{.}\PY{n}{grade\PYZus{}q2b}\PY{p}{(}\PY{n}{q2b\PYZus{}answer}\PY{p}{)} \PY{c+c1}{\PYZsh{} don\PYZsq{}t change this line}
\end{Verbatim}
\end{tcolorbox}

    \begin{Verbatim}[commandchars=\\\{\}]
CORRECT
    \end{Verbatim}


    \newpage

    \hypertarget{normals-and-coins}{%
\subsection{3. Normals and Coins}\label{normals-and-coins}}

Let \(X\) be standard normal. Construct a random variable \(Y\) as
follows:

\begin{itemize}
\tightlist
\item
  Toss a fair coin.
\item
  If the coin lands heads, let \(Y = X\).
\item
  If the coin lands tails, let \(Y = -X\).
\end{itemize}

\textbf{a)} Find the cdf of \(Y\) and hence identify the distribution of
\(Y\).


\begin{align*}
	P(Y \le y) &= P(Y \le y \mid H)P(H) + P(Y \le y \mid  T)P( T) \\
	&= \frac{1}{2} \Phi(y) + \frac{1}{2}P(-X \le y) \\
	&= \frac{1}{2} \Phi(y) + \frac{1}{2} P(X \ge -y) \\
	&= \frac{1}{2} \Phi(y) + \frac{1}{2} (1- \Phi(-y))  \\
	&= \frac{1}{2}\Phi(y) + \frac{1}{2} (1- (1-\Phi(y)), \quad \text{by symmetry of Normal Dist} \\
	&= \frac{1}{2} \Phi(y) + \frac{1}{2}\Phi(y)\\ 
	&= \Phi(y) \\
\end{align*}


Therefore, Y is also standard normal .



\textbf{b)} Find \(E(XY)\) by conditioning on the result of the toss.

\begin{align*}
	E[XY] &= E[E[XY \mid \text{toss}]] \\
	&= E[XY \mid H] P(H) + E[XY \mid T]P(T) \\
	&= \frac{1}{2} E[XX] + \frac{1}{2} E[X(-X)]	\\
	&= \frac{1}{2} E[X^{2}] - \frac{1}{2}E[X^{2}] \\
	&= 0  \\
\end{align*}

\textbf{c)} Are \(X\) and \(Y\) uncorrelated?

\begin{align*}
	\Cov(X, Y) &= E[XY] - E[X] E[Y] \\
	&= 0 - 0(0) \\
	&= 0 \\
\end{align*}

Yes $X$ and $Y$ are uncorrelated.


\textbf{d)} Are \(X\) and \(Y\) independent?

No, $X$ and $Y$ are not independent, because by definition, $Y$  depends on the result of $X$.  

\textbf{e)} Is the joint distribution of \(X\) and \(Y\) bivariate
normal?

No $X$ and $Y$ cannot be bivariate normal. Only bivariate normal implies normal, but normal does not imply bivariate normal. It is only bivariate normal if and only if every linear combination of $X$ and $Y$ is also a normal distribution.   

    \newpage

    \hypertarget{slices-of-a-normal-cake}{%
\subsection{4. Slices of a Normal Cake}\label{slices-of-a-normal-cake}}

This problem needs only the material of a few weeks ago, but it's here
because the ideas and visualization will be helpful in the next
exercise. A former 140 staff member told me how he used this method in
his interview at a major quant firm and surprised the interviewer. (Yes,
he got the job.) Simple, insightful solutions tend to beat tons of
calculus even when the calculus is done correctly.

Let \(X\) and \(Y\) be independent standard normal random variables.

\textbf{a)} Find \(P(X > 0, Y > 0)\).

Yes, it's easy. But get a piece of paper and draw the event on the plane
anyway. Imagine the joint density surface over the plane, and try to
imagine the relevant volume under the joint density surface as a
quadrant-shaped slice of a bell-shaped cake. \textbf{Then use the same
approach for the next two parts.}

\[
P(X > 0, Y > 0) = \frac{1}{4}
\]
\begin{centering}
	
\begin{tikzpicture}
    % 1. Shade the first quadrant (0 to 90 degrees)
    % The 'fill' command fills the area defined by the path
    \fill[blue!20] (0,0) -- (2,0) arc (0:90:2) -- cycle;

    % 2. Draw the full circle
    \draw (0,0) circle (2cm);

    % 3. Optional: Draw the x and y axes
    \draw[->] (-2.5,0) -- (2.5,0) node[right] {$x$};
    \draw[->] (0,-2.5) -- (0,2.5) node[above] {$y$};
\end{tikzpicture}

\end{centering}

\textbf{b)} Find \(P(X > 0, Y > X)\).


\[
P(X > 0, Y > X) = \frac{45}{360} = \frac{1}{8}
\]

\begin{centering}
	
\begin{tikzpicture}
    % 1. Shade the first quadrant (0 to 90 degrees)
    % The 'fill' command fills the area defined by the path
    \fill[blue!20] (0,0) -- (2,0) arc (0:45:2) -- cycle;

    % 2. Draw the full circle
    \draw (0,0) circle (2cm);

    % 3. Optional: Draw the x and y axes
    \draw[->] (-2.5,0) -- (2.5,0) node[right] {$x$};
    \draw[->] (0,-2.5) -- (0,2.5) node[above] {$y$};
\end{tikzpicture}

\end{centering}








\textbf{c)} Find \(P(X > 0, Y > \sqrt{3}X)\).

\[
	P(X > 0, Y > \sqrt{3}X ) = \frac{30}{360} = \frac{1}{12}
\]

\begin{centering}
	
\begin{tikzpicture}
    % 1. Shade the first quadrant (0 to 90 degrees)
    % The 'fill' command fills the area defined by the path
    \fill[blue!20] (0,0) -- (2,0) arc (0:30:2) -- cycle;

    % 2. Draw the full circle
    \draw (0,0) circle (2cm);

    % 3. Optional: Draw the x and y axes
    \draw[->] (-2.5,0) -- (2.5,0) node[right] {$x$};
    \draw[->] (0,-2.5) -- (0,2.5) node[above] {$y$};
\end{tikzpicture}

\end{centering}




    \newpage

    \hypertarget{heights-of-mothers-and-daughters}{%
\subsection{5. Heights of Mothers and
Daughters}\label{heights-of-mothers-and-daughters}}

The heights of a population of mother-daughter pairs have a bivariate
normal distribution with correlation 0.5.

\textbf{a)} Of the mothers on the 90th percentile of mothers' heights,
what proportion have daughters who are taller than the 90th percentile
of daughters' heights?

$\rho = 0.5$ 

Let $X$ be the height of mothers, with a normal distribution. Let $Y$ be the height of daughters. 

Let $Z_{X}$ be the height of mothers with a standard normal distribution, and let $Z_{Y}$ be the height of daughters with a standard normal distribution. $Z_{X} = \frac{X-\mu _{x}}{\sigma_{X}}, Z_{Y} = \rho Z_{X} + \sqrt{1-\rho^{2}} Z$ 

Let $Z$ be a standard normal variable.

\begin{align*}
	P(Z_{Y} \ge Z_{X} \mid Z_{X} = \Phi^{-1}(0.9)) &= P(\rho Z_{X} + \sqrt{1-\rho^{2}}Z  \ge Z_{X} \mid Z_{X} = \Phi^{-1}(0.9)) \\
				    &= P(\sqrt{1-\rho^{2}}Z \ge (1-\rho)Z_{X} \mid Z_{X} = \Phi^{-1}(0.9))  \\
				    &= P(Z \ge \frac{(1-\rho)Z_{X}}{\sqrt{1-\rho^{2}}} \mid Z_{X} = \Phi^{-1}(0.9) ) \\
				    &= P(Z \ge \frac{(1-\rho) \Phi^{-1}(0.9)}{\sqrt{1-\rho^{2}}}) \\
				    &= 1 - \Phi\left( \frac{(1-\rho)\Phi^{-1}(0.9)}{\sqrt{1-\rho^{2}} } \right)  \\
				    &= 0.23 \\
\end{align*}

\textbf{b)} In what proportion of mother-daughter pairs are both women
taller than average? (This means the mothers are taller than the average
mother and the daughters are taller than the average daughter.)

{[}Hint: Express standard bivariate normal variables in terms of two
independent standard normal variables, and then apply the ``slices of a
normal cake'' method.{]}

\begin{align*}
	P(Z_{X} \ge 0, Z_{Y} \ge 0) &= P(Z_{X} \ge 0, \rho Z_{X} + \sqrt{1-\rho^{2}} Z \ge 0) \\
				    &= P(Z_{X} \ge 0, Z \ge -\frac{\rho}{\sqrt{1-\rho^{2}}} Z_{X}) \\
\end{align*}


\begin{align*}
	\alpha = \arctan\left(\frac{\rho}{\sqrt{1-\rho^{2}}} \right)  &=  0.52 \\
	\theta = \alpha + \frac{\pi}{2}
\end{align*}

\begin{align*}
P(Z_{X} \ge 0, Z_{Y}\ge 0) &= \frac{\theta}{2\pi} \\
	&= 0.33 \\
\end{align*}


    \newpage

    \hypertarget{submission-instructions}{%
\subsection{Submission Instructions}\label{submission-instructions}}

Many assignments throughout the course will have a written portion and a
code portion. Please follow the directions below to properly submit both
portions.

\hypertarget{written-portion}{%
\subsubsection{Written Portion}\label{written-portion}}

\begin{itemize}
\tightlist
\item
  Scan all the pages into a PDF. You can use any scanner or a phone
  using applications such as CamScanner. Please \textbf{DO NOT} simply
  take pictures using your phone.
\item
  Please start a new page for each question. If you have already written
  multiple questions on the same page, you can crop the image in
  CamScanner or fold your page over (the old-fashioned way). This helps
  expedite grading.
\item
  It is your responsibility to check that all the work on all the
  scanned pages is legible.
\item
  If you used \(\LaTeX\) to do the written portions, you do not need to
  do any scanning; you can just download the whole notebook as a PDF via
  LaTeX.
\end{itemize}

\hypertarget{code-portion}{%
\subsubsection{Code Portion}\label{code-portion}}

\begin{itemize}
\tightlist
\item
  Save your notebook using
  \texttt{File\ \textgreater{}\ Save\ and\ Checkpoint}.
\item
  Generate a PDF file using
  \texttt{File\ \textgreater{}\ Download\ As\ \textgreater{}\ PDF\ via\ LaTeX}.
  This might take a few seconds and will automatically download a PDF
  version of this notebook.

  \begin{itemize}
  \tightlist
  \item
    If you have issues, please post a follow-up on the general Homework
    12 Ed thread.
  \end{itemize}
\end{itemize}

\hypertarget{submitting}{%
\subsubsection{Submitting}\label{submitting}}

\begin{itemize}
\tightlist
\item
  Combine the PDFs from the written and code portions into one PDF.
  \href{https://smallpdf.com/merge-pdf}{Here} is a useful tool for doing
  so.
\item
  Submit the assignment to Homework 12 on Gradescope.
\item
  \textbf{Make sure to assign each page of your pdf to the correct
  question.}
\item
  \textbf{It is your responsibility to verify that all of your work
  shows up in your final PDF submission.}
\end{itemize}

If you are having difficulties scanning, uploading, or submitting your
work, please read the
\href{https://edstem.org/us/courses/94054/discussion/7533569}{Ed Thread}
on this topic and post a follow-up on the general Homework 12 Ed thread.


    % Add a bibliography block to the postdoc
    
    
    
\end{document}
